\section{Introduction}
\label{sec:intro}

A world model (WM) \citep{ha2018recurrent} constructs the world observations as states and predicts future states based on given actions. Modern world models are trained with massive data \citep{liu2024world,cui2023universal}, demonstrating successful prediction, generation, and planning capabilities. However, existing world models do not directly generalize to structured data, primarily graphs, that are ubiquitous in science \cite{jin2018junction,you2018graph} and industry \cite{ying2018graph,you2022roland} and can be further enriched with multi-modal information \cite{ektefaie2023multimodal}.
% It applies a unified model \citep{liu2024world,cui2023universal} to multiple applications with massive training data, which enhances its capabilities in prediction, generation, and planning. 
% In today's digital era, structured data enriched with multi-modal information is pervasive and essential for a wide range of applications. 
% Although WM has achieved promising results in many applications with unstructured data, its exploration of structured data remains limited. 
Therefore, our paper aims to raise attention to this pressing research question: 
% \textit{Given the importance of managing both unstructured and structured data in diverse applications, how to develop a WM capable of processing these data types simultaneously and generalizing across a broad range of tasks?}
\textit{Can we extend a WM to handle graph-structured data across a broad range of tasks?}

Existing WMs mainly focus on unstructured data. For example, iVideoGPT \citep{wu2024ivideogpt} and Genie \citep{bruce2024genie} are successful world models over video data. However, the relations and structures in the data are rarely explored in these works. Although some works \citep{zhang2021world,zhu2022value} attempt to model structured data in WM using graphs, they have focused solely on planning problems in a specific domain.
% , limiting WM's generalization capabilities across a broader range of tasks. 
In recent years, researchers have also explored the concept of the Graph Foundation Model (GFM)~\citep{liu2023one,chen2024llaga}. However, these methods are confined to predefined graph learning tasks, which cannot easily extend to: (1) multi-modal input data including images and text, (2) diverse tasks beyond standard graph prediction tasks, and (3) data without explicit structure, \ie, standard unstructured data.
% which utilizes a large volume of unstructured training data to build a unified model for addressing multiple tasks.
% limits their potential to become a  WM with strong generalizability.

To address these challenges, we propose the Graph World Model (GWM) that embeds the capabilities of the graph into the WM, which models the current state as a graph and the action as a node (see Table~\ref{Tab:GWM_intro} for comparison with existing methods).
Various tasks can be expressed as action nodes; for example, in a graph prediction task, predicting the label of a given node/edge/subgraph leads to \textit{intended} action nodes that link relevant nodes in the state graph, \ie, target nodes, to the action node; in a retrieval-augmented generation (RAG) task, we can also represent a user query as an \textit{unintended} action node that links to target nodes in the state graph via embedding similarities.
% The nodes of the state graph that are related to the action node are regarded as target nodes. 

To build GWMs, we first introduce a simplified token-based GWM (GWM-T), which integrates multi-modal data like image, table, and text into text modality and represents them as nodes in a graph state. We further develop a token-level message-passing algorithm that aggregates the neighbor information to update the text representation of the state node. Finally, the target nodes on the state graph and prompted action nodes will be fed into multi-modal decoders such as LLMs and Stable Diffusion \cite{rombach2022high}. Despite its simplicity, GWM-T sometimes suffers from high token costs and limited context length. Inspired by latent diffusion models, which introduced modeling in latent space rather than directly on pixels like diffusion to enhance model performance and efficiency, we further develop an embedding-based GWM (GWM-E). GWM-E first employs modality-specific encoders to process different modalities into node embeddings. Then it utilizes embedding-level message
passing to update the node embedding. Finally, the multi-modal information in target state nodes is consolidated through a multi-hop projector before passing them to the decoders. 

We conduct extensive experiments on 6 tasks from diverse domains, including world prediction (multi-modal generation and matching, recommendation, graph prediction), world generation (multi-agent collaboration, retrieval-augmented generation), and world optimization (planning and optimization), with both proposed GWM variants and domain-specific baselines. Results show that (1) GWMs generalize across domains, as the same GWM outperforms or matches domain-specific baselines' performance, (2) graph information matters in GWM, as GWMs benefit from multi-hop graph information, and (3) GWMs demonstrate strong zero-shot/few-shot capabilities on unseen new tasks.

% \todo{Add experiment result summary}


% Extensive experiments on seven representative tasks validated that (1) GWM is able to match the performance of domain-specific methods; (2) GWM benefits from multi-hop graphs; (3) GWM boosts zero-shot/few-shot performance.


% \begin{table}[t]
%     \caption{\textbf{Comparison with existing methods from three perspectives: task type, data structure, and token cost.} Compared with traditional WM and graph FM, GWM can tackle multi-domain tasks and apply to both structured and unstructured data. Additionally, GWM-E is more token-efficient than GWM-T.} %%不要提所有的，提模范典范，examplar 
%     \label{Tab:GWM_intro}
%     \centering
%     \resizebox{0.5\textwidth}{!}{
%     \begin{tabular}{lccc}
%         \toprule
%         \textbf{Method} & \textbf{Task Type} & \textbf{Data Structure} & \textbf{Token Cost} \\
%         \midrule
%         Traditional WMs & Mostly single domain  & Unstructured & - \\
%         Graph FMs  & Graph domain  & Structured & - \\ \midrule
%         \rowcolor{cyan!10} GWM-T  & Multiple domains & Both & High \\
        
%         \rowcolor{cyan!10} GWM-E & Multiple domains & Both & Low \\
%         \bottomrule
%         \end{tabular}}
% \end{table}



\begin{table}[t]
    \setlength\tabcolsep{2pt}
    \caption{\textbf{Comparison with existing representative works from three perspectives: task type, data structure, and model type.} Compared with existing WM and GFM, GWM can tackle multi-domain tasks and be applied to both structured and unstructured data. } 
    \label{Tab:GWM_intro}
    \centering
    \resizebox{0.48\textwidth}{!}{
    \begin{tabular}{lccc}
        \toprule
        \textbf{Method} & \textbf{Task Type} & \textbf{Data Structure} & \textbf{Model Type} \\
        \midrule
         Genie \citep{bruce2024genie} & Video generation  & Unstructured & WM \\
         $L^3P$ \citep{zhang2021world} & Planning  & Structured & WM \\
         BioBridge \citep{wangbiobridge}  & Biomedical domain  & Structured & GFM \\ 
         LLAGA \citep{chen2024llaga}  & Graph domain  & Structured & GFM \\ \midrule
        \rowcolor{cyan!10} GWM-T  & Multiple domains & Both & GWM \\
        \rowcolor{cyan!10} GWM-E & Multiple domains & Both & GWM \\
        \bottomrule
        \end{tabular}}
    % \vspace{-5mm}
\end{table}




